\section{Certified Differential Testing}\label{sec:difftest}
Differential testing runs two implementations of one specification on the
same inputs and reports every input on which they
disagree~\cite{mckeeman1998differential,yang2011csmith}. Its weakness is the
oracle problem: a disagreement says that one side is wrong but not which,
and an agreement says nothing when both sides share a bug or when the inputs
never reach the code that matters. For proof verifiers both problems are
acute: two hand-written verifiers can share a missing check, and inputs
produced by honest provers or by random mutation almost never reach a
verifier's soundness-critical checks, because a mutated proof fails its
first hash comparison. A certified library addresses both.

\paragraph{Setting}
A deployed system implements a proof system with a prover $\Pd$ and a
verifier $\Vd$, following a published specification. We instantiate the
certified library at the system's group, encodings and hash, which gives a
certified verifier $\Vc$ and a certified prover $\Pc$, and compare
$\Vd(x)$ with $\Vc(x)$ on inputs $x$, and both with the verdict we know by
construction. $\Vc$ is not an arbitrary second opinion: it consists of a
certified core and a trusted boundary, and only the boundary is
unverified. The boundary is the \emph{decoder}, which turns the system's
serialised data into group elements and exponents and checks subgroup
membership and ranges; the \emph{challenge closures}, which compute the
Fiat--Shamir challenges with the system's own hash encoding and pass them
to the certified verifier as oracles; and the \emph{wiring}, which decides
which ciphertexts, keys and auxiliary data each proof is about. The core
is the verification function with its theorems. Hence a disagreement has
exactly three possible causes: a bug in $\Vd$, a bug in the boundary of
$\Vc$, or an ambiguity in the specification; Section~\ref{sec:lessons}
shows that all three occur.

\paragraph{Sources of inputs}
We combine eight sources (Table~\ref{tab:sources}). Honest proofs from
$\Pd$ (S1) test the completeness of $\Vc$ and, since we also compare the
challenges both sides derive, our closures against the deployed hash;
honest proofs from $\Pc$ (S2) test $\Vd$ on inputs the deployed prover
never produces, since $\Pc$ takes its randomness as an argument and runs
on all-zero or all-one randomness, the identity permutation, or identity
and repeated ciphertexts. Field mutations (S3) replace one occurrence of
one field by a value of the same type; structural and encoding mutations
(S4) change shapes and encodings (a vector one element short, a
non-minimal or negative integer encoding, an exponent $\ge q$, an integer
outside the group). Mutations of non-hashed fields break a single
verification equation and test each response check in isolation. For
S5, since a mutated proof is rejected at its first hash comparison and
never exercises the algebraic checks, we run $\Pc$ with honest randomness
on a statement and a witness that do not match (a ``permutation'' that
maps two outputs to one input, a re-encryption exponent off by one, a
plaintext changed in one output): every challenge is derived from the
proof's own transcript, so a verifier under test evaluates every
algebraic check on a false statement. S7 tampers with complete signed
election records, re-signed with the development keys the systems
publish; S8 is coverage-guided fuzzing~\cite{jazzer} of the deployed
side's decoders, encoders and hash encodings alone.

\begin{table}[t]
  \caption{Sources of test inputs. ``Accept'' and ``reject'' are the verdict
  both verifiers must give; ``compare'' means only agreement is required.}
  \label{tab:sources}
  \centering\footnotesize
  \begin{tabularx}{\columnwidth}{@{}l>{\raggedright\arraybackslash}Xl@{}}
    \toprule
    Source & Produced by & Verdict \\
    \midrule
    S1 honest proofs & $\Pd$ & accept \\
    S2 honest, degenerate & $\Pc$ (randomness $0$/$1$, identity permutation, trivial ciphertexts) & accept \\
    S3 field mutations & mutate one field by type & reject\,/\,compare \\
    S4 structure, encodings & lengths, shapes, ranges, encodings & reject\,/\,compare \\
    S5 false statements & $\Pc$ on a wrong witness & reject \\
    S6 single-check forgeries & $\Pc$ with one deviation, + control & reject\,/\,accept \\
    S7 protocol tampering & edit and re-sign a whole election & reject \\
    S8 coverage-guided fuzzing & Jazzer~\cite{jazzer} on $\Vd$ alone & no crash, no accept \\
    \bottomrule
  \end{tabularx}
\end{table}

\paragraph{Single-check forgeries (S6)}
The checks of a verifier fall into two classes. \emph{Response checks}
relate responses to commitments; since responses are not hashed, mutating
one breaks only the checks that mention it (S3). \emph{Bound checks}
compare a hashed value with a constant or with the statement, e.g.\
$c_{d,m+1} = 1$ in the zero argument of Bayer--Groth or $E_m = C$ in its
multi-exponentiation argument. Mutating the hashed value changes the
challenge and breaks every equation at once, so no mutation and no random
input can tell whether a verifier implements a given bound check; yet
these are the checks whose omission makes a verifier unsound. For each
bound check $b$ we build, with an independent implementation of the
specification's prover, a \emph{forgery}: the prover run on a false
statement with one deviation chosen so that every check except $b$ holds
(the zero statement sums to $e \ne 0$ and the prover commits honestly to
the diagonal term $e$; the product chain starts at $a_0 + e$; the
commitment opens to $b + e$; the ciphertext is multiplied by
$\mathrm{Enc}(g^e;\sigma)$ and the discrepancy absorbed in $E_m$ or in
$b_m = e$), and a \emph{control}, the same construction with $e = 0$. A
verifier that accepts the control and rejects the forgery demonstrably
implements $b$, and $b$ is demonstrably necessary, since without it the
verifier would accept a false statement; when the verifier reports which
check failed, we confirm that it names $b$ and nothing else.
Table~\ref{tab:bgchecks} lists the seven bound checks of the Swiss Post
Bayer--Groth verifier.

\begin{table}[t]
  \caption{The bound checks of the Bayer--Groth verifier of the Swiss Post
  specification~\cite{swisspost_cryptoprimitives_spec}, and the deviation
  that isolates each.}
  \label{tab:bgchecks}
  \centering\footnotesize\setlength{\tabcolsep}{4pt}
  \begin{tabular}{@{}lll@{}}
    \toprule
    Sub-argument & Bound check & False statement \\
    \midrule
    Zero (Alg.~9.23) & $c_{d,m+1}=1$ & $\sum_i a_i\star b_{i-1}=e$ \\
    Single value product (9.26) & $\tilde b_0=\tilde a_0$ & chain starts at $a_0+e$ \\
     & $\tilde b_{n-1}=x\,b$ & product is $b+e$ \\
    Hadamard (9.21) & $c_{B_0}=c_{A_0}$ & chain starts at $a_0+e$ \\
     & $c_{B_{m-1}}=c_b$ & $c_b$ opens to $b+e$ \\
    Multi-exponentiation (9.16) & $E_m=C$ & $C\cdot\mathrm{Enc}(g^{e};\sigma)$ \\
     & $c_{B_m}=1$ & idem, absorbed by $b_m=e$ \\
    \bottomrule
  \end{tabular}
\end{table}

\paragraph{Engineering}
The deployed verifiers are called in-process through their own APIs and
given exactly the inputs $\Vc$ receives; each certified run is a separate
process on its own input file, so a decoder failure rejects one case and
not a batch. Every external system is pinned to a commit or release with
its toolchain, every random choice of our harnesses is seeded, and where a
deployed prover draws its own randomness its outputs are stored as
fixtures, so that every verdict table of Section~\ref{sec:casestudies} is
reproduced from stored cases by one script per system. \TODO{Swiss Post
harness: seed the extracted prover's generator (currently
\texttt{/dev/urandom}) and store the Java-generated cases as fixtures.}
